\section{Ordinary products in homotopy fixed points}\label{sec:ordinary-products}
We prove Lemma~\ref{lem:lower-bar-product} by constructing relative
products on the ordinary bar tower. The construction supplies cup
products and differential derivations independently of convergence;
separated convergence identifies the final product with the product
on actual homotopy groups.

\begin{proof}[Proof of Lemma~\ref{lem:lower-bar-product}]
Use the homogeneous free bar construction $EQ_\bullet$, with
$EQ_q=Q^{q+1}$, and let $EQ^{(s)}$ be its $s$-skeleton.  The full
totalization tower of the group-action resolution is
\[
 V_s=F_Q((EQ^{(s)})_+,R),\qquad
 \operatorname{holim}_sV_s=F_Q(EQ_+,R)=R^{hQ}.
\]
Here $F_Q$ denotes the derived spectrum of equivariant maps.  The
skeletal filtration gives the usual homogeneous cochains and hence the
displayed $E_2$ page.

For the prismatic geometry below, compare Basterra--Mandell
\cite[proof of Theorem~7.1, pp.~23--26]{basterra-mandell2013}.
Their bar diagonal for a (partial) non-unital $E_1$ algebra over a base
ring spectrum $H$ uses the
same half-cut regions and inverse barycentric coordinates. Here we
construct the equivariant relative maps for $EQ$ and identify the
resulting $E_\infty$ product with the associated-graded product on actual
homotopy groups under the lemma's stated filtration hypothesis.
For barycentric coordinates $t=(t_0,\ldots,t_q)\in\Delta^q$, set
$c_{-1}=0$ and $c_i=t_0+\cdots+t_i$.  Define two barycentric vectors by
\begin{equation}\label{eq:lower-half-cut}
 \begin{split}
 \ell_i(t)&=2\bigl(\min(c_i,\tfrac12)-\min(c_{i-1},\tfrac12)\bigr),\\
 \rho_i(t)&=2\bigl(\max(c_i,\tfrac12)-\max(c_{i-1},\tfrac12)\bigr).
 \end{split}
\end{equation}
They are nonnegative and each has sum one.  Equivalently, their $i$th
coordinates are twice the lengths of the intersections of
$[c_{i-1},c_i]$ with $[0,\tfrac12]$ and $[\tfrac12,1]$.  The formulas
are continuous, including when a cumulative sum equals $\tfrac12$.

For an order-preserving map $\theta:[q]\to[q']$, let
$(\theta_*t)_j=\sum_{\theta(i)=j}t_i$.  Each interval belonging to
$\theta_*t$ is the union of the consecutive intervals belonging to the
corresponding fiber of $\theta$; an empty fiber gives an interval of
length zero.  Additivity of intersection lengths gives
\[
 \ell(\theta_*t)=\theta_*\ell(t),\qquad
 \rho(\theta_*t)=\theta_*\rho(t).
\]
The realization relation is
$[\theta^*\sigma,t]=[\sigma,\theta_*t]$ for $\sigma\in EQ_{q'}$.
Thus for every bar simplex $\sigma\in EQ_q$ the formula
\[
 \delta[\sigma,t]=([\sigma,\ell(t)],[\sigma,\rho(t)])
\]
is compatible with all faces and degeneracies and defines a continuous
map $\delta:EQ\to EQ\times EQ$.  It is $Q$-equivariant because the
action changes $\sigma$ and leaves the barycentric coordinates fixed.
Moreover
\[
 [\sigma,t]\longmapsto
 \bigl([\sigma,(1-\tau)t+\tau\ell(t)],
       [\sigma,(1-\tau)t+\tau\rho(t)]\bigr),
 \qquad 0\leq\tau\leq1,
\]
is compatible with the same identifications and is equivariant.  It is
a homotopy from the ordinary diagonal to $\delta$.

Choose $k$ with $c_{k-1}\leq\tfrac12\leq c_k$.  The support of
$\ell(t)$ is contained in the front face $[0,\ldots,k]$, and the support
of $\rho(t)$ is contained in the back face $[k,\ldots,q]$.  These faces
have dimensions $k$ and $q-k$.  Degenerate faces can only lower those
dimensions.  Consequently
\[
 \delta(EQ^{(s)})\subset
 \bigcup_{i+j\leq s}EQ^{(i)}\times EQ^{(j)}.
\]
This also identifies the cellular chain map.  On the region
$c_{k-1}\leq\tfrac12\leq c_k$, the map
$t\mapsto(\ell(t),\rho(t))$ is a homeomorphism onto
$\Delta^k\times\Delta^{q-k}$, with inverse
\[
 t_i=\ell_i/2\ (i<k),\qquad
 t_k=(\ell_k+\rho_k)/2,\qquad
 t_i=\rho_i/2\ (i>k).
\]
In the oriented coordinates $t_1,\ldots,t_q$ on the source and
$\ell_1,\ldots,\ell_k,\rho_{k+1},\ldots,\rho_q$ on the product, its
Jacobian determinant is $2^q>0$.  The regions meet on their boundaries.
The image of the cellular fundamental chain is therefore
\[
 [\sigma]\longmapsto
 \sum_{k=0}^q
 [\sigma|_{[0,\ldots,k]}]\otimes
 [\sigma|_{[k,\ldots,q]}],
\]
with degenerate terms omitted in normalized chains.  This is the
Alexander--Whitney diagonal, now obtained from the explicit map.

Write $EQ/EQ^{(s)}$ for the based cofiber of
$(EQ^{(s)})_+\to EQ_+$, and set $EQ^{(-1)}=\varnothing$.
The skeletal inclusions are $Q$-CW cofibrations, so these quotients
represent their homotopy cofibers.  For $a,b\geq0$, the skeletal
containment above makes the composite of $\delta$ with the quotient to
$(EQ/EQ^{(a-1)})\wedge(EQ/EQ^{(b-1)})$ constant on
$EQ^{(a+b-1)}$: if $i+j<a+b$, then $i<a$ or $j<b$.  We get a
$Q$-equivariant relative map
\[
 EQ/EQ^{(a+b-1)}\longrightarrow
 (EQ/EQ^{(a-1)})\wedge(EQ/EQ^{(b-1)}).
\]
Multiplication on $R$ and precomposition give
\[
 \begin{split}
 F_Q(EQ/EQ^{(a-1)},R)\wedge F_Q(EQ/EQ^{(b-1)},R)
 \longrightarrow F_Q(EQ/EQ^{(a+b-1)},R).
 \end{split}
\]
These maps are compatible as $a$ and $b$ vary, since they are induced
by the same $\delta$ and the nested skeletal quotients.

For completeness, the associativity needed for powers also holds within
this relative filtration.  Partition $[0,1]$ into any $m$ successive
intervals of positive lengths.  For the $j$th output vector, divide
the lengths of its intersections with $[c_{i-1},c_i]$ by the length of
the $j$th interval.  The preceding naturality and equivariance proof
applies verbatim.  If $k_j$ is an index whose coordinate interval
contains the $j$th cut, the supports lie in the successive faces
$[0,\ldots,k_1]$, $[k_1,\ldots,k_2]$, \ldots,
$[k_{m-1},\ldots,q]$.  Their dimensions sum to $q$; boundary cases
can only reduce the sum.  Thus the same construction gives $m$-fold
relative maps.  Composing cut maps is exactly refinement of the
partition.  The positive interval lengths form a convex simplex, so
the partitions for two parenthesizations are joined by varying those
lengths.  This gives a homotopy through maps with the same relative
skeletal property.  Thus the relative products are associative up to
filtered homotopy, as required for their iterated products.

Put $I_s=F_Q(EQ/EQ^{(s-1)},R)$ for $s\geq0$.  The quotient cofiber
sequence for two successive skeleta gives
\[
 \operatorname{cofib}(I_{s+1}\to I_s)
 \simeq F_Q(EQ^{(s)}/EQ^{(s-1)},R)
 \simeq \operatorname{fib}(V_s\to V_{s-1}).
\]
The fiber sequences $I_s\to R^{hQ}\to V_{s-1}$ form a diagram with
the identity on $R^{hQ}$ and the tower transition on the right.
The stable fiber diagram identifies the connecting maps of the
displayed cofiber sequences with those of the $V_s$ tower, with the
usual rotation signs.  Thus the paired relative spectral sequence is
the ordinary tower spectral sequence with the same filtration index.
The relative products pair these sequences and their exact couples.
On cellular cochains the displayed chain formula is the ordinary cup
product; in inhomogeneous coordinates its value is
\[
 (v\smile w)(g_1,\ldots,g_{a+b})
 =v(g_1,\ldots,g_a)\,(g_1\cdots g_a)
       w(g_{a+1},\ldots,g_{a+b}),
\]
with the usual graded signs.  The cellular boundary formula for a
product, together with the internal Koszul sign, gives the derivation
rule in the paired exact couples and hence on every spectral-sequence
page.  This $E_2$ product is the ordinary cup product used by the
ordinary-to-Tate comparison in Section~\ref{sec:lower-exponent}.
Since each subsequent product
is induced on homology, the products agree on every later page as well.

The fiber sequence for $I_a$ identifies its image in $\pi_*R^{hQ}$ with
$\ker(\pi_*R^{hQ}\to\pi_*V_{a-1})$. After
forgetting the relative conditions, the product just constructed is
homotopic to the actual product on $F_Q(EQ_+,R)$ by the equivariant
homotopy from $\delta$ to the diagonal.  Thus these relative products
multiply the actual filtration groups.  Passing to successive relative
quotients is exactly the product in the paired spectral sequence.
When the stated convergence identifies those quotients with
$E_\infty$, this is the associated-graded homotopy product as stated.
\end{proof}
